\lesson{5}{Nov 17 2021 Wed (16:29:10)}{Molecular Structure}{Unit 3}

\subsubsection*{Electrostatic Forces}

\bf{Valence shell electron pair repulsion (VSEPR)} theory is used to predict the geometry and shape of a molecule, which plays an important role in molecular interactions. This theory is based on the concept that valence electron pairs, whether bonded or nonbonded (lone pairs), repel one another. Lone pairs and bonded pairs will spread out and position themselves around the central atom so they are as far away from each other as possible.

The electron geometry (also called VSEPR shape or VSEPR geometry) of a molecule is determined by the number of \bf{Bonded Electrons} and \bf{Lone Pair Electrons} around the \bf{Central Atom}. \bf{Molecular Shape}, also called the \bf{Molecule Structure}, is a general shape of the molecule that considers the repulsive and attractive forces between electron pairs. More often than not, the \bf{VSEPR} shape and molecule structure are different for a molecule. Two things to remember:

\begin{itemize}
    \item \bf{VSEPR} geometry includes unshared electron pairs, if any, on the central atom that determines the geometry.
    \item The molecule structure is the atoms only and the shape the atoms make when bonded together.
\end{itemize}

Electron pairs exist in many positions, or domains, around a central atom. The number of domains is the sum of bonded atoms and lone pairs of electrons surrounding the central atom.

\subsubsection*{Polarity of Molecules}

\bf{Molecular Polarity} is dependent on the difference in \bf{Electronegativity} between atoms in a compound and the molecular asymmetry of the compound's structure. Polarity also influences other physical properties, such as melting and boiling points, solubility, and the interactions between molecules.

\bf{Nonpolar} molecules have an even distribution of charge due to all \bf{Nonpolar Bonds}, or \bf{Symmetrical Polar Bonds} that "cancel out". \bf{Polar Molecules} have a partial negative end and a partial positive end due to asymmetric arrangement of polar bonds.

\newpage
